\section{Preliminary}\label{Section:Preliminary}
\subsection{Basic notions}
Let us fix our notations. We denote by $\bm{0}_n$ the $n$-vector $(0,0,\ldots,0)$.

For $n=0,1,2,\ldots,$ we set the Pochhammer symbol as
\begin{gather*}
(a)_n:=
\begin{cases}
1,& n=0,\\
a(a+1)\cdots(a+n-1), &n\geq1,
\end{cases}
\end{gather*}
and the $q$-shift factorial is defined by
\begin{gather*}
(a;q)_n:=\prod_{j=0}^{n-1}\big(1-aq^j\big),\qquad (a;q)_\infty:=\prod_{j=0}^\infty \big(1-aq^j\big).
\end{gather*}
For any $a\in\mathbb{C}$, the infinite product $(a;q)_\infty$ is convergent. We set
\begin{gather*} (a_1,a_2,\ldots,a_r;q)_{n}:=\prod\limits_{j=1}^r (a_j;q)_n \qquad \text{for $n=0,1,2,\ldots$ or $n=\infty$}.
\end{gather*}

The theta function with the base $q$ is defined by
\begin{gather*}
\theta_q(x):=\sum_{n\in\mathbb{Z}}q^{\frac{n(n-1)}{2}}x^n,
\end{gather*}
which is holomorphic on $\mathbb{C}^*=\mathbb{C}\setminus\{0\}$. The following properties hold:
\begin{gather}
\theta_q(x) =(q,-x,-q/x;q)_\infty,\label{theta_tpf}\\
\theta_q(q^nx) =x^{-n}q^{-\frac{n(n-1)}{2}}\theta_q(x),\qquad n\in\mathbb{Z},\label{theta_q-diff}\\
\theta_q(qx) =\frac{\theta_q(x)}{x}=\theta_q\left(\frac{1}{x}\right)\label{theta_inverse}.
\end{gather}
For $\lambda\in\mathbb{C}^*$, we set the $q$-spiral $[\lambda;q]:=\lambda q^\mathbb{Z}=\{\lambda q^n ;\,n\in\mathbb{Z}\}$. From the equality~\eqref{theta_tpf}, we see
\begin{gather*}
\theta_q\left(-\frac{\lambda}{x}\right)=0\quad \Leftrightarrow \quad x\in[\lambda;q].
\end{gather*}
We remark that the $q$-spiral $[\lambda;q]$ is identified with an element of $\mathbb{C}^*/q^\mathbb{Z}$.

The $q$-gamma function $\Gamma_q(x)$ is defined by
\begin{gather*}
\Gamma_q(x):=\frac{(q;q)_\infty}{(q^x;q)_\infty}(1-q)^{1-x}.
\end{gather*}
When $q$ is a real number and satisfies $0<q<1$, the limit of $\Gamma_q(x)$ as $q\to1$ gives the gamma function (cf.\ Gasper--Rahman {\cite[Section~1.10]{Gasper&Rahman2004}})
\begin{gather}\label{qgamma_reduction}
\lim_{q\to1}\Gamma_q(x)=\Gamma(x).
\end{gather}

\subsection[Review of $q$-Borel summability]{Review of $\boldsymbol{q}$-Borel summability}\label{Review}
We review the theory of $q$-Borel summability (cf.~Zhang \cite{Zhang2002, Zhang2005} and Dreyfus--Eloy~\cite{DreyfusEloy2016}).

Let $\mathbb{C}[[x]]$ be the ring of formal power series of $x$ and $\hat{f}(x)=\sum\limits_{n\ge0}a_nx^n\in\mathbb{C}[[x]]$. For $\hat{f}(x)$, we define its $q$-Borel transform $\hat{\mathcal{B}}_{q;1}\colon \mathbb{C}[[x]]\to\mathbb{C}[[\xi]]$ by
\begin{gather*}
(\hat{\mathcal{B}}_{q;1}\hat{f})(\xi):=\sum_{n\ge0}a_nq^{\frac{n(n-1)}{2}}\xi^n.
\end{gather*}

We denote $\mathcal{M}(\mathbb{C}^*,0)$ the field of functions that are meromorphic on some punctured neigh\-borhood of~0 in $\mathbb{C}^*$. More explicitly $f\in\mathcal{M}(\mathbb{C}^*,0)$ if and only if there exists $V$, an open neighborhood of 0, such that $f$ is analytic on $V\setminus\{0\}$.

Let $[\lambda;q]\in\mathbb{C}^*/q^\mathbb{Z}$ and $f\in\mathcal{M}(\mathbb{C}^*,0)$. It is said that $f$ belongs to $\mathbb{H}_{q;1}^{[\lambda;q]}$ if there exist a positive constant $\varepsilon$ and a domain $\Omega\subset\mathbb{C}$ such that:
\begin{itemize}\itemsep=0pt
\item $\displaystyle\bigcup_{m\in\mathbb{Z}}\left\{x\in\mathbb{C}^* ;\,\left|x-\lambda q^m\right|<\varepsilon\left|q^m\lambda\right|\right\}\subset\Omega$.
\item The function $f$ can be continued to an analytic function on $\Omega$ with $q$-exponential growth at infinity, which means that there exist positive constants $L$ and $M$ such that for any $x\in\Omega\setminus\{0\}$, the following holds:
\begin{gather*}
|f(x)|<L\theta_{|q|}(M|x|).
\end{gather*}
We note that this estimate does not depend on the choice of $\lambda$.
\end{itemize}

Let $[\lambda;q]\in\mathbb{C}^*/q^{\mathbb{Z}}$. For $f\in\mathbb{H}_{q;1}^{[\lambda;q]}$, we define its $q$-Laplace transform $\mathcal{L}_{q;1}^{[\lambda;q]}\colon \mathbb{H}_{q;1}^{[\lambda;q]}\to\mathcal{M}(\mathbb{C}^*,0)$ by
\begin{gather*}
\big(\mathcal{L}_{q;1}^{[\lambda;q]}f\big)(x):=\frac{1}{1-q}\int_{0}^{\lambda\infty}\frac{f(\xi)}{\theta_q\big(\frac{\xi}{x}\big)}\frac{d_q\xi}{\xi}=\sum_{n\in\mathbb{Z}}\frac{f(\lambda q^n)}{\theta_q\big(\frac{\lambda q^n}{x}\big)}.
\end{gather*}
Here, this transformation is given in Jackson integral (cf.\ Gaspar--Rahman \cite[Section~1.10]{Gasper&Rahman2004}). For sufficient small $|x|$, the function $\big(\mathcal{L}_{q;1}^{[\lambda;q]}f\big)(x)$ has poles of order at most one that are contained in the $q$-spiral $[-\lambda;q]$. We remark that the well-definedness of $q$-Laplace transform can be shown in the same way as Th\'eor\`eme~1.3.2 of Zhang~\cite{Zhang2002}. If the function $\hat{f}(x)\in\mathbb{C}[[x]]$ satisfies $\big(\hat{\mathcal{B}}_{q;1}\hat{f}\big)\in\mathbb{H}_{q;1}^{[\lambda;q]}$, we call that $\hat{f}(x)$ is $[\lambda;q]$-summable and $\big(\mathcal{L}_{q;1}^{[\lambda;q]}\circ\hat{\mathcal{B}}_{q;1}\hat{f}\big)(x)$ is the $[\lambda;q]$-sum of~$\hat{f}(x)$.

We give some fundamental properties of $q$-Borel--Laplace transform without proofs.
\begin{Proposition}[Dreyfus--Eloy {\cite[Section~1]{DreyfusEloy2016}}] Let $[\lambda;q]\in\mathbb{C}^*/q^\mathbb{Z}$ and $f$ be an analytic function. Then, the followings hold:
\begin{itemize}\itemsep=0pt
\item $\big(\hat{\mathcal{B}}_{q;1}f\big)\in\mathbb{H}_{q;1}^{[\lambda;q]}$,
\item $\big(\mathcal{L}_{q;1}^{[\lambda;q]}\circ\hat{\mathcal{B}}_{q;1}f\big)=f$.
\end{itemize}
\end{Proposition}
\begin{Proposition}\label{prop_borel2} Let $[\lambda;q]\in\mathbb{C}/q^\mathbb{Z}$. For $\hat{f}(x)\in\mathbb{C}[[x]]$ and $g\in\mathbb{H}_{q;1}^{[\lambda;q]}$, the followings hold.
\begin{itemize}\itemsep=0pt
\item $\big(\hat{\mathcal{B}}_{q;1}\big(x^m\sigma_q^n\hat{f}\big)\big)=q^{\frac{m(m-1)}{2}}\xi^m\sigma_q^{n+m}\big(\hat{\mathcal{B}}_{q;1}\hat{f}\big)$,
\item $\big(\mathcal{L}_{q;1}^{[\lambda;q]}(\xi^m\sigma_q^ng)\big)=q^{-\frac{m(m-1)}{2}}x^m\sigma_q^{n-m}\big(\mathcal{L}_{q;1}^{[\lambda;q]}g\big)$.
\end{itemize}
\end{Proposition}
The first equality is shown by direct calculation. The second equality can be proved by using the property (2.2) of theta function. From this proposition, we have the following important corollary.
\begin{Corollary}\label{cor_q-diff}
Let $\hat{f}(x)\in\mathbb{C}[[x]]$ be a formal solution of a~linear $q$-difference equation. If $\hat{f}(x)$ is $[\lambda;q]$-summable, its $[\lambda;q]$-sum satisfies the same $q$-difference equation with $\hat{f}(x)$.
\end{Corollary}

At the end of this section, we give asymptotic properties of the $q$-Borel--Laplace transform.

Let $[\lambda;q]\in\mathbb{C}^*/q^\mathbb{Z}$, $\hat{f}(x)=\sum\limits_{n\ge0}a_nx^n\in\mathbb{C}[[x]]$ and $f\in\mathcal{M}(\mathbb{C}^*,0)$. Then we define
\begin{gather*}
f\sim_1^{[\lambda;q]}\hat{f},
\end{gather*}
if for any positive numbers $\varepsilon$ and $R$, there exist positive constants $C$ and $K$ such that for any $N\ge1$ and
\begin{gather*}
x\in\left\{x\in\mathbb{C}^* ;\,|x|<R\right\}\setminus\bigcup_{m\in\mathbb{Z}}\left\{x\in\mathbb{C}^* ;\,|x+\lambda q^m|<\varepsilon |q^m\lambda |\right\},
\end{gather*}
we have
\begin{gather*}
\left|f(x)-\sum_{n=0}^{N-1}a_nx^n\right|\le LM^N|q|^{-\frac{N(N-1)}{2}}|x|^N.
\end{gather*}
\begin{Proposition}[Dreyfus--Eloy {\cite[Section~1]{DreyfusEloy2016}}]\label{qBorel_AE}
Let $[\lambda;q]\in\mathbb{C}^*/q^\mathbb{Z}$ and $\hat{f}\in\mathbb{C}[[x]]$ with $\big(\hat{\mathcal{B}}_{q;1}\hat{f}\big)\in\mathbb{H}_{q;1}^{[\lambda;q]}$. Then $\big(\mathcal{L}_{q;1}^{[\lambda;q]}\circ\hat{\mathcal{B}}_{q;1}\hat{f}\big)\sim_{1}^{[\lambda;q]}\hat{f}$.
\end{Proposition}

\section{Main results}\label{Section:Main Results}
In this section, we give our main results which are about the resummation of the basic hypergeometric series
\begin{gather}\label{qHGS2}
\hat{f}(x)={}_r\varphi_s(\bm{a};\bm{b};q,x)={}_r\varphi_s\left( \begin{matrix}\bm{a}\\ \bm{b} \end{matrix};q,x\right).
\end{gather}
We remark that $\hat{f}(x)$ is the formal power series solution of the linear $q$-difference equation \eqref{qHGEq}. Our results are stated as follows.
\begin{Theorem}\label{Theorem:Main1}
We put $k=r-s-1$ and $p=q^k$. Assuming $a_i/a_j\notin q^{\mathbb{Z}}$, $i\neq j$, basic hypergeometric series \eqref{qHGS2} is $[\lambda;p]$-summable for any $\lambda\in\mathbb{C}^*\setminus\big[(-1)^{k};q\big]$ and its $[\lambda;p]$-sum ${}_rf_s(\bm{a};\bm{b};\lambda;q,x)=\big(\mathcal{L}_{p;1}^{[\lambda:p]}\circ\mathcal{\hat{B}}_{p;1}\hat{f}\big)(x)$ is given by
\begin{gather}
{}_rf_s(\bm{a};\bm{b};\lambda;q,x) =\frac{(a_2,\ldots,a_r,b_1/a_1,\ldots,b_{s}/a_1;q)_\infty}{(b_1,\ldots,b_{s},a_2/a_1,\ldots,a_r/a_1;q)_\infty}\frac{\theta_{p}\big(pa_1^kx/\lambda\big)}{\theta_{p}(px/\lambda)}
\frac{\theta_q\big((-1)^{1-k}a_1\lambda\big)}{\theta_q\big((-1)^{1-k}\lambda\big)}\nonumber\\
\hphantom{{}_rf_s(\bm{a};\bm{b};\lambda;q,x) =}{}
\times{}_r\varphi_{r-1}\left(\begin{matrix}a_1,a_1q/b_1,\ldots,a_1q/b_s,\bm{0}_k\\a_1q/a_2,\ldots,a_1q/a_r
\end{matrix} ;q,\frac{qb_1\cdots b_s}{a_1\cdots a_rx}\right)\nonumber\\
\hphantom{{}_rf_s(\bm{a};\bm{b};\lambda;q,x) =}{}
+\idem(a_1;a_2,\ldots,a_r),\label{Main1}
\end{gather}
where $x\in\mathbb{C}^*\setminus[-\lambda;p]$ and $|qb_1b_2\cdots b_s/a_1\cdots a_rx|<1$. Here the symbol ``$\,\idem(a_1;a_2,\ldots,a_r)$'' means the sum of $r-1$ terms which are obtained by interchanging $a_1$ with each $a_2,\ldots,a_r$ in the preceding expression.
The function ${}_rf_s(\bm{a};\bm{b};\lambda;q,x)$ has simple poles on the $p$-spiral $[-\lambda;p]$.
\end{Theorem}
From Corollary~\ref{cor_q-diff} and the well-definedness of $q$-Laplace transform, we see that the $[\lambda;p]$-sum ${}_rf_s(\bm{a};\bm{b};\lambda;q,x)=\big(\mathcal{L}_{p;1}^{[\lambda:p]}\circ\mathcal{\hat{B}}_{p;1}\hat{f}\big)(x)$ itself becomes the meromorphic solution of~\eqref{qHGEq} on some punctured neighborhood of the origin in $\mathbb{C}^*$. The expression~\eqref{Main1} gives the analytic continuation of ${}_rf_s(\bm{a};\bm{b};\lambda;q,x)$ to the neighborhood of infinity.
 In addition, this shows a $q$-analog of the Stokes phenomenon. Indeed, an actual solution~\eqref{Main1} has~\eqref{qHGS2} as its asymptotic expansion (see Proposition~\ref{qBorel_AE}) and its value changes depending on the choice of $[\lambda;p]\in\mathbb{C}^*/p^{\mathbb{Z}}$.
Using the fundamental system of solutions around infinity \eqref{qHGEq_Sol}, the $[\lambda;p]$-sum \eqref{Main1} can be expressed as
\begin{gather*}
{}_rf_s(\bm{a};\bm{b};\lambda;q,x)=\sum_{j=1}^r M_j(x,\lambda)y_j^{(\infty)}(x),
\end{gather*}
where $M_j(x,\lambda)=C_jT_j(x,\lambda)$ and
\begin{gather*}
C_j=\frac{(b_1/a_j,b_2/a_j,\ldots,b_s/a_j;q)_\infty}{(b_1,b_2,\ldots,b_s;q)_\infty}\prod_{1\le i\le r,\, i\neq j}\frac{(a_i;q)_\infty}{(a_i/a_j;q)_\infty},\\
T_j(x,\lambda)=\frac{\theta_{p}\big(pa_j^kx/\lambda\big)}{\theta_{p}(px/\lambda)}\frac{\theta_q\big((-1)^{1-k}a_j\lambda\big)}{\theta_q\big((-1)^{1-k}\lambda\big)}\frac{\theta_q(-x)}{\theta_q(-a_jx)}\label{T_j}.
\end{gather*}
The coefficients $M_j$ are called $q$-Stokes coefficients. Unlike the case of differential equation, the values of coefficients $M_j$ change continuously dependent on $\lambda\in\mathbb{C}^*\setminus\big[(-1)^k;q\big]$.
